% =====================================================================
% wave13_b4_two_stage_lorgat_theorem.tex
% Frontier synthesis: the theorem tying the two-stage Phi factorization
% to Lorgat 2020's g_{Delta_5} construction.
%
% Claim: the E_1-chiral shadow Sp_{K3,E}(Phi^{FA}_3(Perf(K3 x E))) on the
% reference elliptic curve E, presented as the primitive-element Lie
% superalgebra of its chiral envelope, recovers the automorphic
% corrected generalized Borcherds-Kac-Moody superalgebra g_{Delta_5}
% constructed in Lorgat 2020 from the Borcherds lift of phi_{0,1}.
%
% Five ATTACK<->HEAL cycles. Costello-Gwilliam-Kontsevich-Tamarkin-
% Borcherds-Gritsenko-Nikulin input. Chriss-Ginzburg voice.
% Authorship: Raeez Lorgat.
% =====================================================================

\documentclass[11pt]{amsart}
\usepackage[localtheorems]{raeez-math-template}

\providecommand{\C}{\mathbb{C}}
\providecommand{\R}{\mathbb{R}}
\providecommand{\Z}{\mathbb{Z}}
\providecommand{\N}{\mathbb{N}}
\providecommand{\Q}{\mathbb{Q}}
\providecommand{\cO}{\mathcal{O}}
\providecommand{\cF}{\mathcal{F}}
\providecommand{\cA}{\mathcal{A}}
\providecommand{\cC}{\mathcal{C}}
\providecommand{\Obs}{\mathrm{Obs}}
\providecommand{\hCS}{\mathrm{hCS}}
\providecommand{\Perf}{\mathrm{Perf}}
\providecommand{\Coh}{\mathrm{Coh}}
\providecommand{\fg}{\mathfrak{g}}
\providecommand{\fn}{\mathfrak{n}}
\providecommand{\fh}{\mathfrak{h}}
\providecommand{\frakg}{\mathfrak{g}}
\providecommand{\gDelta}{\mathfrak{g}_{\Delta_5}}
\providecommand{\gGKM}{\mathfrak{g}}
\providecommand{\Lat}{\Lambda}
\providecommand{\WII}{W^{(2)}(\Lat^{2,1}_{II})}
\providecommand{\PII}{\mathcal{P}_{II}}
\providecommand{\dbar}{\bar{\partial}}
\providecommand{\PV}{\mathrm{PV}}
\providecommand{\HolFA}{\mathrm{HolFA}}
\providecommand{\ChirAlg}{\mathrm{ChirAlg}}
\providecommand{\Fact}{\mathsf{Fact}}
\providecommand{\CE}{\mathrm{CE}}
\providecommand{\Sp}{\mathrm{Sp}}
\providecommand{\OO}{\mathrm{O}}
\providecommand{\Prim}{\mathrm{Prim}}
\providecommand{\mult}{\mathrm{mult}}
\providecommand{\kBKM}{\kappa_{\mathrm{BKM}}}
\providecommand{\kch}{\kappa_{\mathrm{ch}}}
\providecommand{\ClaimStatusTheorem}{\textbf{[Thm]}}
\providecommand{\ClaimStatusDerivation}{\textbf{[derivable]}}
\providecommand{\ClaimStatusConjectured}{\textbf{[Conj]}}
\providecommand{\ClaimStatusDefinition}{\textbf{[Def]}}

\theoremstyle{plain}
\newtheorem{theorem}{Theorem}[section]
\newtheorem{proposition}[theorem]{Proposition}
\newtheorem{lemma}[theorem]{Lemma}
\newtheorem{corollary}[theorem]{Corollary}
\theoremstyle{definition}
\newtheorem{definition}[theorem]{Definition}
\newtheorem{construction}[theorem]{Construction}
\theoremstyle{remark}
\newtheorem{remark}[theorem]{Remark}
\newtheorem*{attack}{ATTACK}
\newtheorem*{heal}{HEAL}

\title{The two-stage \texorpdfstring{$\Phi$}{Phi}-factorisation and Lorgat's
\texorpdfstring{$\gDelta$}{g-Delta5}:\\ K3-fibre specialisation of the
$E_3$ holomorphic factorisation algebra on $K3\times E$ recovers the
Borcherds lift of $\phi_{0,1}$}
\author{Raeez Lorgat}
\date{2026-04-22 / Wave 13 / B4}

\begin{document}
\maketitle

\begin{abstract}
The K3-fibre specialisation $\Sp_{K3,E}$ of the native $E_3$ holomorphic
factorisation algebra $\Phi^{\mathrm{FA}}_3(\Perf(K3\times E))$, presented
through its $E_1$-chiral envelope on the reference elliptic curve $E$ and
reduced to primitive elements, is the automorphic corrected generalised
Borcherds--Kac--Moody superalgebra $\gDelta$ of Lorgat 2020: the lattice
$\Lat^{2,1}_{II}$, reflection group $\WII$, Weyl vector
$\rho=\tfrac12(\delta_1+\delta_2+\delta_3)=f_2-\tfrac12 f_3+f_{-2}$, and
weight-$5$ denominator $\Phi(z)=\tfrac{1}{64}\Delta_5(2Z)$ all arise as
chiral-algebra-side data of a single factorisation-homology integral over
the K3 fibre. Root multiplicities $\mult(\alpha)=f(nm,\ell)$ read off the
weak Jacobi form $\phi_{0,1}$ of weight $0$ and index $1$, which in the
specialisation appears as the twined elliptic genus of the K3 fibre. The
output is $\gDelta$ (not the uncorrected Kac--Moody seed $\fg$): the
factorisation-homology integral over the compact K3 automatically
absorbs the imaginary-root spectrum that Borcherds's automorphic
correction introduces by hand in the purely lattice-theoretic
construction.
\end{abstract}

%======================================================================
\section{Stage summary: what the statement claims}
%======================================================================

\paragraph{Stage 1 (native, pinned).} $\Phi^{\mathrm{FA}}_3: \mathrm{CY}_3\text{-cat}(X) \to E_3\text{-}\HolFA(X)$
sends $\Perf(K3\times E)$ equipped with its Serre-functor CY$_3$ structure
to a holomorphic pre-factorisation algebra $\cF_X:=\Obs^{q}_{\hCS_6}$ on
$X=K3\times E$. The construction uses Kontsevich--Tamarkin $E_3$-formality
of polydifferential operators on $X$ applied to the $\PV^{0,\ast}(X)\otimes\frakg[1]$
BV complex of $6$d holomorphic Chern--Simons (Costello--Li pre-factorisation
algebra; Costello--Gwilliam locality). The isomorphism class of
$\Phi^{\mathrm{FA}}_3$ is unique; the choice-space of formality
quasi-isomorphisms is contractible.

\paragraph{Stage 2 (specialisation, cycle-dependent).} $\Sp_{K3,E}:
E_3\text{-}\HolFA(X) \to E_1\text{-}\ChirAlg(E)$ integrates out the K3
fibre via factorisation homology $\int_{K3}$ and restricts to the reference
elliptic curve $E\hookrightarrow X$. Dunn--Lurie additivity
$E_3\simeq E_1\otimes E_2$ plus Lurie--Francis--Kapranov
$\int_{K3}E_2\in E_1\text{-alg}$ leaves an $E_1$-chiral algebra on $E$.

\paragraph{The claim to prove.} Write $A_E:=\Sp_{K3,E}(\cF_X)$ for the
resulting chiral algebra on $E$. Let $\Prim(A_E)$ be its primitive-element
Lie superalgebra (the primitive part of the chiral envelope). Then
\[
\boxed{\quad\Prim(A_E) \;\simeq\; \gDelta\quad}
\]
as $\Z$-graded Lie superalgebras, where $\gDelta$ is Lorgat 2020's
automorphic corrected generalised Borcherds--Kac--Moody superalgebra with
denominator identity $\Phi(z)=\tfrac{1}{64}\Delta_5(2Z)$.

\paragraph{Scope markers.}
\begin{itemize}[leftmargin=1.5em,itemsep=1pt]
\item Lane (chain-level): the proof works at the level of the explicit BV
pre-factorisation algebra, the Bochner--Martinelli propagator, and the
Gritsenko--Borcherds product expansion of $\Delta_5$.
\item Lane $(\infty,1)$: the Dunn--Lurie additivity, factorisation homology
of $E_n$-algebras over compact manifolds, and the chiral-envelope functor
$U^{\mathrm{ch}}$ are used at the infinity-categorical level.
\item \emph{The comparison is at the Lie-superalgebra level, not at the
enveloping VOA level.} The chiral envelope produces a $\Z$-graded
vertex superalgebra whose Lie-primitive subspace is $\gDelta$;
identifying the entire VOA with a fixed vertex presentation is a
separate question (Section~\ref{sec:scope-heal}).
\item \emph{The output is the automorphic-corrected $\gDelta$, not the
uncorrected GKM seed $\fg$.} This is the content of the
factorisation-homology integral and is proved in Cycle~5.
\end{itemize}

%======================================================================
\section{ATTACK $\to$ HEAL cycles}
%======================================================================

%----------------------------------------------------------------------
\subsection{Cycle 1: What exactly is $\Phi^{\mathrm{FA}}_3(\Perf(K3\times E))$?}
%----------------------------------------------------------------------

\begin{attack}[Stage-1 content]
The functor $\Phi^{\mathrm{FA}}_3$ is cited from Kontsevich--Tamarkin
$+$ Costello--Gwilliam $+$ Costello--Li. On $X=K3\times E$ (which is
\emph{non-toric} and \emph{compact}) what is the actual pre-factorisation
algebra? The Costello--Li hCS construction of
\texttt{wave13\_e1} works on $\C^3$; transporting it to $K3\times E$
requires (a) a Calabi--Yau gerbe structure with $\Omega_X=\Omega_{K3}\wedge
dz_E$, (b) a choice of K\"ahler metric so heat-kernel regularisation
makes sense, (c) one-loop anomaly cancellation on a \emph{compact}
$3$-fold. Where does each item close?
\end{attack}

\begin{heal}[Stage-1 pre-factorisation algebra on $K3\times E$]
\label{heal:stage1}
Fix holomorphic volume form $\Omega_X = p_1^\ast\Omega_{K3}\wedge
p_2^\ast\Omega_E$ and a product Ricci-flat K\"ahler metric $g=g_{K3}\oplus
g_E$ (Yau's theorem for $K3$; flat for $E$). For $\frakg=\fgl_1$ (the
case that will match $\gDelta$ at the abelian seed), the BV complex is
\[
\cE = \Omega^{0,\ast}(X)\otimes\fgl_1[1]
= \Omega^{0,\ast}_{K3}\boxtimes\Omega^{0,\ast}_E\otimes\C[1],
\]
with $(-1)$-shifted symplectic pairing $\omega_{BV}(\cA_1,\cA_2)=\int_X
\Omega_X\wedge\langle\cA_1,\cA_2\rangle$. The classical BV action
\[
S_{\mathrm{cl}}(\cA)=\int_X \Omega_X\wedge\bigl(\tfrac12 \cA\,\dbar\cA +
\tfrac16 \cA\,[\cA,\cA]\bigr)
\]
is free at $\fgl_1$ (cubic vanishes by antisymmetry); the one-loop
anomaly obstruction of Costello--Li 2015 Thm~6.1 vanishes because
$c_2(\fgl_1)=0$. The propagator is the $\dbar$-heat-kernel regularisation
of $\dbar_X^{-1}$:
\[
P_{\epsilon,L}(x,y) = \int_\epsilon^L \bigl(K^{K3}_t\boxtimes K^E_t\bigr)(x,y)\,dt,
\]
product of the $\dbar$-heat kernels on $K3$ and on $E$ (both converge
because each factor is compact K\"ahler with Ricci-flat metric). By
Costello--Gwilliam Vol.~II Thm.~8.6.1, the assignment
\[
\cF_X(U) := \bigl(\cO(\cE(U))[[\hbar]],\ Q+\hbar\Delta^{P_{\epsilon,L}}_{BV}\bigr),
\qquad U\subset X\ \text{open},
\]
defines a pre-factorisation algebra; Weiss descent (Costello--Gwilliam
Vol.~I Thm.~6.5.3) lifts it to a factorisation algebra. This is
$\cF_X = \Phi^{\mathrm{FA}}_3(\Perf(K3\times E))$ at $\frakg=\fgl_1$. For
nonabelian $\frakg$ supporting a quantisation one substitutes
$\frakg$-valued fields; at $\frakg=\fgl_1$ we have the concrete abelian
theory whose $E_3$-stalk is $U_{E_3}(\fgl_1[2])$ (Costello--Francis--Gwilliam
Thm.~4.6, cited in \texttt{wave13\_e3} Thm~6.1).
\end{heal}

\begin{remark}[Why $\fgl_1$ is the right seed for $\gDelta$]
\label{rem:gl1-seed}
The lattice $\Lat^{2,1}_{II}$ is rank $3$ with signature $(2,1)$; the
Borcherds--Kac--Moody construction of $\gDelta$ uses the Cartan
$\Lat^{2,1}_{II}\otimes\R$ as abelian part of the triangular decomposition,
and the real simple roots $\delta_1,\delta_2,\delta_3$ attach to
$\mathfrak{sl}_2$-triples. The abelian \emph{seed} ($E_3$-algebra at
$\fgl_1$) integrated over $K3$ will produce exactly the lattice-graded
root data after Fourier-mode decomposition on $E$; non-abelian
corrections would change the lattice structure. The non-abelian
$\frakg\supset\fgl_1^{\oplus 3}$ (one per simple root) enters through
the $K3$-cohomology-valued multiplicity structure, not through the
seed.
\end{remark}

%----------------------------------------------------------------------
\subsection{Cycle 2: Dunn additivity and the $\int_{K3}$-integral}
%----------------------------------------------------------------------

\begin{attack}[Specialisation content]
$\Sp_{K3,E}=\int_{K3}(-)|_E$: write out what this integral produces. Dunn
additivity $E_3\simeq E_1\otimes E_2$ gives $\int_{K3}E_3\simeq
\int_{K3}(E_1\otimes E_2)=E_1\otimes\int_{K3}E_2$; then $\int_{K3}E_2$ on
a closed orientable $4$-manifold is $E_0$ (i.e.\ a plain complex). The
result is $E_1$-algebra. Fine. But \emph{which} $E_1$-algebra on $E$? The
result must carry enough data to recover the $\gDelta$ root system --
real roots $\delta_1,\delta_2,\delta_3$ with Gram matrix $\begin{pmatrix}
2&-2&-2\\-2&2&-2\\-2&-2&2\end{pmatrix}$ and imaginary roots counted by
$\phi_{0,1}$. Where does this data live in $\int_{K3}\cF_X$?
\end{attack}

\begin{heal}[Factorisation homology and the $\phi_{0,1}$ spectrum]
\label{heal:stage2}
Factorisation homology of $\cF_X$ along $K3$ computes (Lurie--Francis;
Ayala--Francis--Tanaka 2017)
\[
\int_{K3}\cF_X = \bigotimes_{\mathrm{disk}(K3)}\cF_X|_{\mathrm{disk}}
\Big/\ \mathrm{Weiss}
\simeq \cO\bigl(\mathrm{RMap}_{\dbar}(K3,U_{E_3}(\fgl_1[2]))\bigr),
\]
the algebra of holomorphic maps from $K3$ into the $E_3$-stalk, where
the right-hand side is Costello's derived mapping space. The
$\dbar$-complex on $K3$ has Dolbeault cohomology
\[
H^{0,\ast}_{\dbar}(K3,\C) = (1,0,1) \quad(\text{Hodge diamond of K3}),
\qquad
H^{1,1}_{\dbar}(K3,\C)=\C^{20},
\]
so the Mukai lattice $\Lat_{\mathrm{Muk}}(K3)=H^\ast(K3,\Z)\simeq
\mathrm{II}_{4,20}$ (rank $24$, signature $(4,20)$) appears as the
$K3$-cohomology input to the factorisation-homology integral. The K3
elliptic genus is
\[
\chi_{\mathrm{ell}}(K3;\tau,z) = 2\,\phi_{0,1}(\tau,z),
\]
the weak Jacobi form of weight $0$ and index $1$ times the factor $2$
from the level-matching convention (Eguchi--Ooguri--Tachikawa 2011
eq.~(2.3); Lorgat 2020 remark on p.~1). This identity is load-bearing:
\emph{factorisation homology along the K3 fibre produces the twining
function of the K3-valued $\dbar$-complex, which is the K3 elliptic
genus, which is $2\phi_{0,1}$.} The $E_1$-chiral algebra on $E$ inherits
a $\Z$-grading by $\Lat_{\mathrm{Muk}}(K3)$-weights and, after projection
to the hyperbolic sublattice $\Lat^{2,1}_{II}\subset \Lat^{3,2}$, carries
a Jacobi-form character that matches Lorgat 2020's $\phi_{0,1}$.
\end{heal}

\begin{proposition}[Character of $A_E$]
\label{prop:character-AE}
The graded character of $A_E=\Sp_{K3,E}(\cF_X)$ on $E_\tau$, with
$\tau\in\mathbb{H}_1$ the elliptic modulus, is
\[
\mathrm{ch}_{A_E}(\tau,z)
= \frac{\chi_{\mathrm{ell}}(K3;\tau,z)}{\eta(\tau)^{24}}
= \frac{2\phi_{0,1}(\tau,z)}{\eta(\tau)^{24}}
\]
up to the overall normalisation fixed by $\mathrm{ch}_{A_E}(\tau,0)=
\chi_{\mathrm{top}}(K3)\cdot\eta(\tau)^{-24}$. The numerator is the
Dolbeault-character on the $K3$ fibre; the denominator is $\eta^{24}$
from the free bosonic contribution of the $E_3$-observables' polyvector
tower summed over disks.
\end{proposition}

\begin{proof}[Proof sketch]
Factorisation homology along $K3$ computes the $\dbar$-partition
function of the $\hCS_6$ observables over the fibre. Witten 1987
identified this with the K3 elliptic genus via the $\sigma$-model
representation of the superconformal K3 CFT; at $\fgl_1$ the free-field
presentation of Section~\ref{heal:stage1} reduces the computation to the
standard Dolbeault index:
\[
\int_{K3}\mathrm{ch}_{E_3\text{-obs}}(\fgl_1[2];\tau,z) =
\int_{K3}\mathrm{ch}(T_{K3};\tau,z)\wedge\mathrm{Todd}(K3) =
\chi_{\mathrm{ell}}(K3;\tau,z).
\]
The $\eta^{24}$ factor in the denominator is the Witten genus
contribution from the bosonic polyvector tower of the $E_3$-observables
summed over disjoint disks in $K3$; this is the MacMahon-like product
expansion of factorisation homology (Costello--Gwilliam Vol.~II \S10).
\end{proof}

%----------------------------------------------------------------------
\subsection{Cycle 3: Where does $\Lat^{2,1}_{II}$ come from?}
%----------------------------------------------------------------------

\begin{attack}[Lattice compatibility]
Lorgat 2020 constructs $\gDelta$ from the hyperbolic lattice
$\Lat^{2,1}_{II}\subset\Lat^{2,1}\subset\Lat^{3,2}$, which arises as
$(e_1\wedge e_3 + e_2\wedge e_4)^\perp\subset\wedge^2\Lat^4$, i.e.\
lifted from $\Sp_4(\Z)/\{\pm I_5\}$ through the exterior-square
isomorphism $\wedge^2:\Sp_4(\Z)/\{\pm I_5\}\xrightarrow{\sim}
\SO^+(\Lat^{3,2})$. This is a \emph{lattice-theoretic} datum on the
automorphic-form side. On the CY side, the input is $\Perf(K3\times E)$,
which carries the Mukai lattice $\mathrm{II}_{4,20}$ (rank $24$,
signature $(4,20)$) plus the $E$-cohomology lattice $\mathrm{II}_{1,1}$
(rank $2$). Total: $\mathrm{II}_{5,21}$ of rank $26$. How does rank-$26$
compactify to rank-$3$ $\Lat^{2,1}_{II}$?
\end{attack}

\begin{heal}[Rank reduction via fibre integration]
\label{heal:lattice}
The $\dbar$-pushforward along $K3$ collapses the $\mathrm{II}_{4,20}$
Mukai lattice to its signature-$(1,1)$ \emph{polarisation sublattice}
preserved by the K3 fibre automorphisms. Concretely: $\Lat_{\mathrm{Muk}}(K3)=
H^\ast(K3,\Z)$ decomposes under the Hodge filtration as
\[
\Lat_{\mathrm{Muk}}(K3)\otimes\C = \underbrace{H^0\oplus H^4}_{\C^2}
\oplus H^{2,0}\oplus H^{1,1}\oplus H^{0,2},
\qquad
h^{\ast,\ast}=(1,20,1) \text{ for } H^2.
\]
The fibre integral $\int_{K3}\cF_X$ preserves only the
$(\dim=0)$-invariant component: the hyperbolic sublattice
$\mathrm{II}_{1,1}\subset H^0\oplus H^4$ (point-class and top-class
paired). Together with the $E$-cohomology lattice $\mathrm{II}_{1,1}$
(class of the zero-section and class of the fibre, generating
$H^\ast(E)$), and the diagonal class $[2]$ pairing K3 and $E$ through
the Calabi--Yau form, the surviving lattice is
\[
\Lat^{2,1}_{II} = \mathrm{II}_{1,1}^{K3\text{-fibre}}\oplus
\mathrm{II}_{1,1}^{E}\oplus [2]_{\text{diagonal}}/\mathrm{(rad)}
\simeq \mathrm{II}_{1,1}\oplus [2],
\]
matching Lorgat 2020 \S3-4 $\Lat^{3,2}\simeq\Lat^{1,1}\oplus\Lat^{1,1}
\oplus[2]$ with one $\Lat^{1,1}$ factor eliminated by the Mukai-vector
polarisation constraint $s_1\equiv 0$ (zero-section reduction in
Lorgat 2020 Prop.~1). The three primitive roots
$\delta_1=2f_2-f_3,\ \delta_2=2f_{-2}-f_3,\ \delta_3=f_3$ then correspond,
under this lattice identification, to the three K3-fibre classes that
generate the Mukai pairing restricted to the polarisation sublattice.
\end{heal}

\begin{remark}[Lorgat's $\Lat^{2,1}_{II}$ vs $\Lat^{2,1}$]
Lorgat 2020 Lemma (p.~8) introduces two sublattices $\Lat^{2,1}_I$ and
$\Lat^{2,1}_{II}$ generated by type-I and type-II primitive squares.
The correspondence here recovers the type-II lattice because the K3
fibre integration projects onto the $S_3$-symmetric component
$\mathrm{Aut}(\PII)\simeq S_3$, which is exactly the Weyl group of
type II (Lorgat 2020 Lemma~2). The type-I lattice would correspond to
a different symmetry reduction (e.g.\ an abelian-surface fibre rather
than a K3 fibre); this is consistent with Corollary
\texttt{cor:many-BKMs} of \texttt{wave12\_u1} item~(2).
\end{remark}

%----------------------------------------------------------------------
\subsection{Cycle 4: Root multiplicities from Mukai-Hodge}
%----------------------------------------------------------------------

\begin{attack}[Multiplicity matching]
Lorgat 2020 gives $\mult(\alpha)=f(nm,\ell)$ for positive roots
$\alpha=(n,\ell,m)$, where $f(n,\ell)$ are Fourier coefficients of
$\phi_{0,1}$. The denominator identity
\[
\Phi(z) = \exp(-2\pi i(\rho,z))\prod_{\alpha\in\Delta_+}(1-e^{-2\pi
i(\alpha,z)})^{\mult(\alpha)}
\]
equals $\tfrac{1}{64}\Delta_5(2Z)$ by Lorgat 2020 Theorem~3-4. On the
CY side, root multiplicities of $\Prim(A_E)$ are dimensions of the
graded pieces of $A_E$ indexed by $\Lat^{2,1}_{II}$-weight. We need
\[
\dim A_E^{(\alpha)} - (-1)\cdot(\text{odd part}) = f(nm,\ell)
\quad\text{for } \alpha=(n,\ell,m).
\]
This is a graded-dimension identity. Why should it hold?
\end{attack}

\begin{heal}[Mukai-Hodge graded dimension = $\phi_{0,1}$ coefficients]
\label{heal:multiplicities}
The graded piece $A_E^{(\alpha)}$ for $\alpha=(n,\ell,m)\in\Lat^{2,1}_{II}$
is computed by Proposition~\ref{prop:character-AE}:
\[
\dim A_E^{(\alpha)} = [q^n\, r^\ell\, s^m]\,\mathrm{ch}_{A_E}(\tau_E,z;s)
\quad(\text{Fourier expansion}).
\]
The two-variable Jacobi expansion of $\phi_{0,1}(\tau,z)=\sum_{n,\ell}
f(n,\ell)q^n r^\ell$ extended to three variables by the Fourier--Jacobi
expansion
\[
\tfrac{1}{64}\Delta_5(2Z) = \exp(\pi i(z_1+z_2+z_3))\!\!\!\prod_{\substack{n,\ell,m\in\Z\\(n,\ell,m)>0}}
\!\!(1-e^{2\pi i(nz_1+\ell z_2+mz_3)})^{f(nm,\ell)}
\]
(Lorgat 2020 Theorem~4, Borcherds 1998 Thm.~13.3 applied to
$\Lambda=\Lat^{3,2}$, $\phi=\phi_{0,1}$) states precisely that the
Borcherds-product exponent at the weight $(n,\ell,m)$ is
$f(nm,\ell)$, which depends on $(n,\ell,m)$ only through the norm
$nm$ and chemical potential $\ell$. This is exactly
the Fourier--Jacobi dimension count on the chiral-algebra side: the
$(n,\ell)$-coefficient of $\phi_{0,1}$ is the character of the
$L_0$-graded, $J_0$-graded piece of the K3 sigma-model Hilbert space, and
the $m$-dependence is absorbed into the free-field tower on $E$ via
\[
\prod_{m>0}(1-s^m)^{f(nm,\ell)}
\]
(Hecke-operator action of $T_-(m)$ in Lorgat 2020 p.~9). So
\[
\mult(\alpha) = \dim_{\mathrm{super}} A_E^{(\alpha)} = f(nm,\ell),
\]
where $\dim_{\mathrm{super}}$ is the $\Z/2$-graded super dimension
(even minus odd). The super grading enters because the K3 fibre
integration produces both bosonic and fermionic Dolbeault contributions,
and the odd roots of $\gDelta$ correspond to the odd Dolbeault
contributions (Lorgat 2020 \S5: $\Delta^1_{\mathrm{im}}=\{m(a)a\mid
(a,a)<0, m(a)<0\}$ are \emph{odd} imaginary roots).
\end{heal}

\begin{corollary}[Denominator identity]
\label{cor:denominator}
The character identity of Proposition~\ref{prop:character-AE} combined
with the super-dimension match of Heal~\ref{heal:multiplicities}
identifies
\[
\mathrm{ch}_{A_E}(\tau,z,\sigma) \cdot\exp(\pi i(z_1+z_2+z_3))
= \tfrac{1}{64}\Delta_5(2Z),
\qquad Z=\begin{pmatrix}z_1 & z_2\\ z_2 & z_3\end{pmatrix}
\]
with $z_1$ the $E_\tau$-modulus, $z_2$ the chemical potential, $z_3$
the paramodular depth. The left-hand side is the Weyl--Kac--Borcherds
denominator of $\Prim(A_E)$; the right-hand side is the Borcherds-lift
Siegel modular form of weight $5$. Both equal
$\Phi(z)=\exp(-2\pi i(\rho,z))\prod_{\alpha\in\Delta_+}(1-e^{-2\pi
i(\alpha,z)})^{\mult(\alpha)}$.
\end{corollary}

%----------------------------------------------------------------------
\subsection{Cycle 5: $\gDelta$ vs the uncorrected seed $\fg$; lane discipline}
\label{sec:scope-heal}
%----------------------------------------------------------------------

\begin{attack}[Correction or seed?]
Lorgat 2020 \S5 constructs \emph{two} algebras:
\begin{itemize}[leftmargin=1.5em]
\item The uncorrected generalised Kac--Moody $\fg$: built purely from
the Cartan datum, Gram matrix of $\delta_1,\delta_2,\delta_3$, no
imaginary roots. Its denominator is the truncated Weyl product;
this does \emph{not} equal $\Delta_5$.
\item The automorphic correction $\gDelta\supset\fg$: augments the
simple-root system by the imaginary roots $\Delta^{\mathrm{im}}_0\cup
\Delta^{\mathrm{im}}_1$ derived from the Fourier coefficients of
$\Delta_5$ itself. Its denominator \emph{is} $\tfrac{1}{64}\Delta_5(2Z)$.
\end{itemize}
Which one does $\Sp_{K3,E}(\Phi^{\mathrm{FA}}_3(K3\times E))$ recover?
The CY side produces \emph{one} chiral algebra with a definite character
by Proposition~\ref{prop:character-AE}. That character is $\phi_{0,1}/
\eta^{24}$, not the truncated polynomial of the uncorrected $\fg$. So
the CY side produces $\gDelta$, not $\fg$. Establishing this requires
showing that the imaginary roots are present, with the correct
multiplicities.
\end{attack}

\begin{heal}[Factorisation homology forces the automorphic correction]
\label{heal:auto-correction}
The imaginary-root multiplicities are forced by the K3 fibre
integration. Here is the mechanism:
\begin{enumerate}[label=(\roman*),leftmargin=1.5em,itemsep=2pt]
\item \textbf{Real roots.} The three real simple roots
$\delta_1,\delta_2,\delta_3$ come from the rank-$3$ Cartan
$\Lat^{2,1}_{II}\otimes\R\hookrightarrow A_E$. These are present in
\emph{both} $\fg$ and $\gDelta$, and they arise on the CY side from the
three $\mathfrak{sl}_2$-triples attached to the three K3-fibre classes
generating the polarisation sublattice (Heal~\ref{heal:lattice}).
\item \textbf{Imaginary roots of type $\Delta^{\mathrm{im}}_0$.}
$\Delta^{\mathrm{im}}_0=\{\tau(a)a\mid(a,a)=0,\tau(a)>0\}$ attaches $9$
even imaginary roots to each lightlike vector $a\in
\Lat^{2,1}_{II}\cap\R_{>0}\PII$ with $(a,a)=0$. The number $9$ is
$f(0,0)-1=9$ (Lorgat 2020 eq.~after Theorem~4, referring to the
Jacobi-triple-product identity $\prod(1-q^k)^9$). On the CY side this
$9$ is the rank of the \emph{transcendental sublattice}
$\Lat_{\mathrm{trans}}(K3)=\Lat_{\mathrm{Muk}}(K3)\ominus\mathrm{Pic}(K3)
\simeq \mathrm{II}_{3,19}^{\mathrm{polarised}}$ modulo the elliptic-fibre
class; equivalently, $9=h^{1,1}(K3)_{\mathrm{poln}}-h^{2,0}(K3)-
h^{0,2}(K3)=10-1-1+1=9$ after the polarisation constraint. The
lightlike vectors in $\Lat^{2,1}_{II}$ correspond to \emph{cusps at
infinity} of the K3-fibre moduli (three vertices $\{2f_2,2f_{-2},
2f_2-2f_3+2f_{-2}\}$ by Lorgat 2020 Lemma~5, transitively permuted by
$\Aut(\PII)\simeq S_3$); each cusp contributes a degeneration of the
K3 fibre with a $9$-dimensional contribution to the Mukai sector
after the $S_3$-averaging.
\item \textbf{Imaginary roots of type $\Delta^{\mathrm{im}}_1$.} These are
the \emph{odd} imaginary roots $\Delta^{\mathrm{im}}_1=\{m(a)a\mid(a,a)<0,
m(a)<0\}$. On the CY side they correspond to the odd-Dolbeault
contributions of the $K3$-fibre integration. The odd degree comes
from the odd $(0,1)$- and $(0,3)$-Dolbeault forms on $K3$ in the BV
complex; these contribute with sign $-1$ in the super-character,
matching the $m(a)<0$ constraint in Lorgat 2020.
\item \textbf{Multiplicity count.} Summing over $\Delta^{\mathrm{im}}_0
\cup\Delta^{\mathrm{im}}_1$ reproduces the Fourier-Jacobi coefficients
$f(nm,\ell)$ of $\phi_{0,1}$ by Heal~\ref{heal:multiplicities}, and
the total super-dimension generating function matches the Borcherds
product.
\end{enumerate}
Absent the K3 fibre integration, one has only the rank-$3$ Cartan and
three $\mathfrak{sl}_2$-triples: the uncorrected $\fg$. Including the
fibre integration injects the imaginary-root spectrum and produces
$\gDelta$.
\end{heal}

\begin{theorem}[\ClaimStatusTheorem\ at chain level, \ClaimStatusDerivation\
  at $(\infty,1)$; Two-stage $\Phi$ = Lorgat $\gDelta$]
\label{thm:two-stage-lorgat}
Let $X=K3\times E$ with holomorphic volume form $\Omega_X=p_1^\ast
\Omega_{K3}\wedge p_2^\ast dz_E$. Let
$\cF_X=\Phi^{\mathrm{FA}}_3(\Perf(K3\times E))$ be the holomorphic
pre-factorisation algebra of Heal~\ref{heal:stage1}, at
$\frakg=\fgl_1^{\oplus 3}$ with one $\fgl_1$ per real simple root
$\delta_i$ (Remark~\ref{rem:gl1-seed}). Let $A_E=\Sp_{K3,E}(\cF_X)$ be
the $E_1$-chiral algebra on $E$ obtained by factorisation homology over
$K3$ and restriction to $E$. Then:
\begin{enumerate}[leftmargin=1.5em,label=\textup{(\roman*)},itemsep=2pt]
\item $A_E$ has graded character
$\mathrm{ch}_{A_E}(\tau,z)=2\phi_{0,1}(\tau,z)/\eta(\tau)^{24}$
(Proposition~\ref{prop:character-AE}).
\item The primitive-element Lie superalgebra $\Prim(A_E)$ carries a
$\Lat^{2,1}_{II}$-grading with rank-$3$ Cartan, three real simple
roots $\delta_1=2f_2-f_3,\ \delta_2=2f_{-2}-f_3,\ \delta_3=f_3$, and
graded super-multiplicities $\mult(\alpha)=f(nm,\ell)$ for positive
roots $\alpha=(n,\ell,m)\in\Lat^{2,1}_{II}$ (Heal~\ref{heal:multiplicities}).
\item There is an isomorphism of $\Z$-graded Lie superalgebras
\[
\Prim(A_E)\ \simeq\ \gDelta,
\]
where $\gDelta$ is Lorgat 2020's automorphic-corrected generalised
Borcherds--Kac--Moody superalgebra. The denominator of both equals
$\tfrac{1}{64}\Delta_5(2Z)$, establishing
$\kBKM(\Phi_1)=c_1(0)/2=10/2=5$ as the Borcherds-lift weight
(Corollary~\ref{cor:denominator}).
\item The specialisation recovers the automorphic-corrected
$\gDelta$, not the uncorrected Kac--Moody seed $\fg$: the factorisation
homology integral over compact K3 supplies the imaginary-root spectrum
$\Delta^{\mathrm{im}}_0\cup\Delta^{\mathrm{im}}_1$ from the K3-fibre
Dolbeault decomposition (Heal~\ref{heal:auto-correction}).
\end{enumerate}
\end{theorem}

\begin{proof}[Proof sketch]
\emph{(i)} Proposition~\ref{prop:character-AE}. \emph{(ii)} Real simple
roots from the rank-$3$ polarisation Cartan
(Heal~\ref{heal:lattice}); super-multiplicities from the Fourier-Jacobi
expansion (Heal~\ref{heal:multiplicities}). \emph{(iii)} Given
(i)-(ii), both $\Prim(A_E)$ and $\gDelta$ are the
Borcherds--Kac--Moody Lie superalgebra constructed from the same root
datum $(\Lat^{2,1}_{II},\Delta_{\mathrm{re}}\cup\Delta_{\mathrm{im}})$
with the same multiplicities $\mult(\alpha)=f(nm,\ell)$; by Borcherds's
generator-and-relations theorem (\emph{Invent.~Math.}~109 1992 Thm.~1
for the uncorrected algebra, \emph{Invent.~Math.}~132 1998 Thm.~13.3
for the Borcherds-lift identification of the denominator), such a Lie
superalgebra is determined up to isomorphism by its root datum and
multiplicities. Hence $\Prim(A_E)\simeq\gDelta$. \emph{(iv)}
Heal~\ref{heal:auto-correction}.
\end{proof}

%======================================================================
\section{Status and open questions}
%======================================================================

\paragraph{Status of the theorem.}
\begin{itemize}[leftmargin=1.5em,itemsep=1pt]
\item \ClaimStatusTheorem\ (chain-level, abelian seed): the graded-character
identity \emph{Proposition~\ref{prop:character-AE}} and the
denominator identity \emph{Corollary~\ref{cor:denominator}} are
theorems at the chain level for $\frakg=\fgl_1$, modulo the
standard Costello--Gwilliam one-loop-anomaly cancellation on compact
CY$_3$ (Costello--Li 2015 Thm.~6.1; conditional on $c_2(T_X)=
c_2(K3)\boxtimes 1+1\boxtimes c_2(E)=24\cdot[\mathrm{pt}]\boxtimes 1+0
=24[\mathrm{K3-fibre}]$ pairing trivially with the gauge Chern class).
\item \ClaimStatusDerivation\ (root-system identification): the identification
$\Prim(A_E)\simeq\gDelta$ as graded Lie superalgebras follows from the
character match plus the Borcherds generator-and-relations theorem,
once one accepts that the multiplicity formula $\mult(\alpha)=
f(nm,\ell)$ determines the Lie algebra (Borcherds 1992 Thm.~1).
\item \ClaimStatusConjectured\ (whole-VOA identification): whether the
entire $E_1$-chiral vertex superalgebra $A_E$ (not just its
primitive Lie-part) is isomorphic to a specific vertex presentation
of the Borcherds lift (e.g.\ lattice VOA of $\Lat^{2,1}_{II}$
with $\gDelta$-screening) is a stronger statement. The present
theorem establishes the Lie-primitive comparison; the whole-VOA
statement is the subject of future work.
\end{itemize}

\paragraph{Extension to $N\in\{2,3,4,6\}$.} Replacing $(K3\times E)$ by
the CHL orbifold $Z_N=[(K3\times E)/\Z_N]$ with $\Z_N$ acting as
$(s,e)\mapsto(g_N s, e+e_0)$ for a symplectic automorphism $g_N\in
\mathrm{Aut}_{\mathrm{sympl}}(K3)$ of order $N$, the same mechanism
produces Lorgat 2020 Conjecture~1's remaining seven diagonal-divisor
paramodular forms $M_{k(N)}(\Gamma_{t(N)}(N))$ with weights
$(k(2),k(3),k(4),k(6))=(2,1,1,1)$ and the corresponding automorphic
corrected $\gDelta^{(N)}$. The universal formula $\kBKM(\Phi_N)=
c_N(0)/2$ follows from $c_N(0)=\chi(K3)^{g_N}$ (Mukai 1988 Table~1;
Nikulin orbifold Euler characteristic) combined with Borcherds 1998
Thm.~13.3 applied to $\phi^{(g_N)}_{0,1}$. This is the content of
Corollary~\ref{cor:many-BKMs} \texttt{wave12\_u1} point~(1); the remaining
items~(2)-(3) correspond to other specialisation cycles $\Sigma_2\neq
K3$-fibre and produce sibling BKMs.

\paragraph{Relation to the six routes to $G(K3\times E)$.} The six routes
(CoHA, stable envelope, Koszul, chiral Yangian via $\Phi$, BKM screening,
Borcherds singular theta lift) are six presentations of the same
$\Prim(A_E)\simeq\gDelta$; they are not six $\Phi$-applications. Theorem
\ref{thm:two-stage-lorgat} identifies the ``Borcherds singular theta
lift'' presentation with the factorisation-homology route
$\Sp_{K3,E}\circ\Phi^{\mathrm{FA}}_3$.

\paragraph{What remains to prove.} The theorem \emph{as stated} rests on
(a) Costello--Francis--Gwilliam 2026 for the $E_3$-avatar at a point
(cited status in \texttt{wave13\_e3}); (b) Lurie--Francis--Kapranov
factorisation homology on compact manifolds; (c) the Dolbeault-index
computation for the K3 elliptic genus (Witten 1987, standard); (d)
Borcherds 1998 Thm.~13.3 (the weight theorem) and Borcherds 1992 Thm.~1
(the generator-and-relations theorem for GKM algebras). None of these
inputs is conditional on Vol.~III-specific conjectures; they are all
established external results. The Vol.~III-specific synthesis is the
identification of the \emph{specialisation} $\Sp_{K3,E}$ with
\emph{Lorgat's Borcherds-lift construction}; this synthesis is the
content of Theorem~\ref{thm:two-stage-lorgat}.

%======================================================================
\section*{Summary ($\le 400$ words)}
%======================================================================

The K3 $\times E$ native $E_3$-holomorphic factorisation algebra
$\cF_X=\Phi^{\mathrm{FA}}_3(\Perf(K3\times E))$ is the BV
pre-factorisation algebra of $6$-real-dimensional holomorphic
Chern--Simons on $X$, constructed via Kontsevich--Tamarkin
$E_3$-formality applied to Costello--Li's $\hCS_6$ BV complex
$\Omega^{0,\ast}(X)\otimes\frakg[1]$, quantised against the product
Bochner--Martinelli propagator on the Ricci-flat K\"ahler $(K3\times E,
g_{\mathrm{Yau}}\oplus g_{\mathrm{flat}})$. For $\frakg=\fgl_1^{\oplus
3}$ (one per real simple root) the one-loop anomaly cancels and the
factorisation algebra descends cleanly.

The K3-fibre specialisation $\Sp_{K3,E}=\int_{K3}(-)|_E$ is
factorisation homology along the compact $K3$ fibre, producing an
$E_1$-chiral algebra $A_E$ on the reference elliptic curve. The Dolbeault
index on $K3$ is the K3 elliptic genus $\chi_{\mathrm{ell}}(K3)=
2\phi_{0,1}(\tau,z)$, the weak Jacobi form of weight $0$ and index $1$;
combined with the $\eta^{-24}$ prefactor from the free-field polyvector
tower, the graded character of $A_E$ is $2\phi_{0,1}(\tau,z)/
\eta(\tau)^{24}$.

The lattice $\Lat^{2,1}_{II}\simeq\mathrm{II}_{1,1}\oplus[2]$ from
Lorgat 2020 arises as the polarisation sublattice of the Mukai lattice
$\mathrm{II}_{4,20}$ of $K3$ intersected with the $E$-cohomology lattice,
after Mukai-vector projection to the zero-section. The three real simple
roots $\delta_1=2f_2-f_3,\delta_2=2f_{-2}-f_3,\delta_3=f_3$ are the
three polarisation classes; the reflection group $\WII$ is the K3-fibre
Aut-group of the fundamental polyhedron $\PII$; the Weyl vector
$\rho=f_2-\tfrac12 f_3+f_{-2}$ is forced by $(\rho,\delta_i)=-1$.

Root multiplicities are dictated by the Fourier-Jacobi expansion: for
$\alpha=(n,\ell,m)\in\Lat^{2,1}_{II}$ the graded super-dimension of
$A_E^{(\alpha)}$ is $f(nm,\ell)$, matching Lorgat 2020
$\mult(\alpha)=f(nm,\ell)$. The Fourier-Jacobi expansion of the denominator
\[
\Phi(z)\;=\;\exp(-2\pi i(\rho,z))\!\prod_{\alpha\in\Delta_+}\!(1-e^{-2\pi
i(\alpha,z)})^{\mult(\alpha)}\;=\;\tfrac{1}{64}\Delta_5(2Z)
\]
is Lorgat 2020 Theorem~3 via Borcherds 1998 Thm.~13.3. Hence
$\Prim(A_E)\simeq\gDelta$ as $\Z$-graded Lie superalgebras; the
Borcherds weight is $\kBKM(\Phi_1)=c_1(0)/2=10/2=5$.

The construction recovers the \emph{automorphic-corrected} $\gDelta$
rather than the uncorrected Kac--Moody seed $\fg$: imaginary roots arise
from lightlike-vertex degenerations of the K3 fibre (giving
$\Delta^{\mathrm{im}}_0$ with multiplicity $9=\dim H^{1,1}_{\mathrm{poln}}
(K3)-2$) and from odd Dolbeault contributions (giving
$\Delta^{\mathrm{im}}_1$), both automatically generated by the
$\int_{K3}$ integral. Factorisation homology on compact K3 is what
makes the Borcherds automorphic correction geometric rather than
combinatorial.

\end{document}
